\section{Related work and theoretical foundations}
Linear forecasting remains a necessary reference in comparisons with deep temporal models. DLinear and NLinear question whether architectural complexity alone explains forecasting accuracy \cite{dlinear}. PatchTST uses channel-independent patches for Transformer forecasting \cite{patchtst}; TimesNet represents temporal variation in two dimensions \cite{timesnet}; iTransformer places variates in the token dimension \cite{itransformer}; TiDE uses dense temporal encoding \cite{tide}; and TimeMixer combines information across temporal scales \cite{timemixer}. These methods have different benchmark recipes and budgets. We execute a small official-code PatchTST configuration as a feasible modern comparator, without transferring published accuracy claims to our protocol. The other architectures are reviewed rather than reported as executed experiments.

Rubin distinguishes missingness mechanisms by the information on which availability depends \cite{rubinmissing}. The lag-based mask here is a specified simulation based on past observations. Missingness-aware recurrent models use both observations and the measurement process. GRU-D introduces input and hidden-state decay \cite{grud}; our shared-service adaptation is not a reproduction of its clinical multivariate setting. Graph-based forecasting with missing values is addressed by BiTGraph \cite{bitgraph} and GinAR \cite{ginar}. Their historical implementations and results are preserved, including very large errors, but are not newly fitted in every transfer scenario. The two GRU comparators in the new experiment test local information and static neighborhood aggregation, using the recurrent cell of \cite{gru}. The ChronoGraph project combines microservice temporal data and call relationships \cite{chronograph}; a known static graph is treated as an input assumption, not as proof that every edge was available at every historical forecast origin.

Cloud-specific studies distinguish the prediction target and the action taken from it. TempoScale concerns microservice autoscaling \cite{temposcale}, AERO includes adaptive period detection in resource prediction \cite{aero}, and STAMP integrates system state and spatiotemporal dependencies for microservice workload prediction \cite{stamp}. Broader reviews discuss AI applications to microservices and autoscaling methods \cite{aimapping,autoscaling2025}. Our authors' recent systematic review provided candidate references; the primary publications were checked separately before inclusion. Its counts of papers are not counts of experiments in this study. Further work on missing-variable forecasting, dynamic telemetry sampling, adaptive microservice management and cloud anomaly detection provides adjacent context, with preprint or under-review status recorded where applicable \cite{ibn,radar2026,taxonomy2026,ibm2026}. Anomaly-detection scores are not directly comparable with forecasting MASE.

STEAM builds a graph from workload similarity and uses system-state information to create alternative graph views \cite{steamgraph}. This provides a point of comparison for the static call adjacency used here. DeepScaling evaluates a scaling policy that targets stable CPU utilization \cite{deepscaling}. Its unit of evaluation includes control decisions, whereas this study measures forecast error on recorded usage.

Finally, uncertainty about a measured mean difference is distinct from uncertainty about an individual future observation. We report conditional resampling intervals for algorithm comparisons using temporal blocks \cite{blockbootstrap}. We do not construct prediction intervals or claim coverage; conformal prediction introduces additional assumptions and calibration decisions beyond these experiments \cite{conformal}. MASE is used as a scale-normalized error with a training-prefix denominator \cite{mase}; it does not remove differences in target semantics between the two data sources.

